\section{实验数量与方案确定顺序}
\label{app:ledger}
\begin{table}[!htbp]\centering\small
\caption{正式观测窗口及事件文件数量。后期实验为静默段单独保存文件，各阶段的文件分段方式不同。}
\begin{tabularx}{\linewidth}{L{3.5cm}rrY}\toprule
实验 & 观测窗口 & 事件文件 & 输入与条件\\\midrule
首轮M0/M1比较 & 280 & 60 & 1701--1705；20条序列，含恢复探测\\
状态干预 & 90 & 90 & 复用1701--1705\\
前瞻全脑检验 & 111 & 36 & 2701--2703；102个M1与9个M0初态窗口\\
经典特征补测 & 525 & 1026 & 3701--3703；含6个空白窗口\\\midrule
合计 & 1006 & 1212 & 一个连接组，11个种子标识\\\bottomrule
\end{tabularx}\end{table}

用于拟合的60个窗口和回顾性检查的80个窗口均为首轮数据子集。102个新预测与102个M1观测逐点配对。正式总数不包含工程预试验、未完成片段及此前糖/苦味基线。同一输入和训练前缀可能在不同分支中复用，因此窗口数和文件数均不作为独立样本量。

\begingroup\small
\begin{longtable}{L{3.1cm}L{5.5cm}L{5.8cm}}
\caption{各项分析的设计与数据使用顺序。}\label{tab:chronology}\\
\toprule 项目 & 方案确定时可用的数据 & 分析性质\\\midrule\endfirsthead
\toprule 项目 & 方案确定时可用的数据 & 分析性质\\\midrule\endhead
\bottomrule\endfoot
首轮正式协议 & 工程预试验结束后，固定输入、参数和判据。 & 正式仿真前在本地保存方案，未作外部预注册。\\
损失恢复比例F & 首轮部分结果已经读取；后续实验前另行规定双归一化指标。 & 首次使用为探索性分析，后续按各自方案执行。\\
状态干预 & 方案先于正式运行确定；首次旁路状态恢复出现实现错误。 & 修复后完成实现检查，再运行正式比较。\\
代理拟合与检查 & 60个拟合窗口及80个检查窗口的输出均已读取。 & 前者估计参数，后者作回顾性检查。\\
102窗口新预测 & 系数、候选条件和逐点预测先行保存，新全脑输出随后产生。 & 前瞻检验，未以新输出重新拟合。\\
局部响应回升 & 已观察到首个强刺激实现的响应上升。 & 事后检查全部九条新序列，原评分不变。\\
H3部分恢复 & 已读取单轮A基线，尚未读取主多轮及恢复输出；参考文献后固定2s方案。 & 补充实验，2s与10s结果分别按各自判据分析。\\
H4(b)补测 & 原比较结果显示仅一例满足全部条件。 & 结果驱动的0.75s/24次补充实验，阈值保持不变。\\
本报告 & 所有实验与分析均已完成。 & 事后综合及文字修订，无新增仿真或拟合。\\
\end{longtable}\endgroup

\section{补充数值与判据}
\subsection{逐输入实现的响应}
\begin{table}[!htbp]\centering\small\setlength{\tabcolsep}{4pt}
\caption{首轮M1的十条序列。除D外，单位为每300ms窗口的脉冲数。M0在1701--1704中恒为10，在1705中恒为9。}
\begin{tabular}{rrrrrrrr}\toprule
种子 & $T$(s) & $I$ & $E$ & $L$ & $D$(\%) & $P_2$ & $P_{10}$\\\midrule
\inputrows{data/individual.tex}\bottomrule
\end{tabular}\end{table}

首轮五种输入的命令事件数分别为2125、2080、2145、2171和2094，实际感觉脉冲数分别为2100、2054、2109、2134和2065。在同一实现内，两者均逐次稳定。阶段二进一步验证了离散输入接口的映射：只有$\mathrm{tick}>\mathrm{last}+1$的命令事件被接受，相应感觉事件出现在$\mathrm{tick}+1$。这一映射与旧数据及前瞻实验的实际感觉事件逐点一致。

\subsection{资源恢复的解析表达}
静默期间资源按下式恢复：
\begin{equation}
x(t)=1-[1-x(t_0)]\exp[-(t-t_0)/\tau_{\rm rec}].
\end{equation}
取$\tau_{\rm rec}=2$s时，等待2s或10s后，资源亏损分别保留约37\%或0.67\%。这是状态方程的解析结果。连续均值输入下，近似稳态为$1/(1+Uf\tau_{\rm rec})$；150Hz时为0.25。该近似对应持续输入，不直接描述200ms间歇刺激的周期状态，也未包含资源至aBN1响应的非线性转换。

\clearpage
\subsection{操作性判据}
\begingroup\small
\begin{longtable}{L{2.7cm}L{11.7cm}}
\caption{主要判据及其先决条件。阈值沿用相应实验方案。}\\
\toprule 项目 & 定义\\\midrule\endfirsthead
\toprule 项目 & 定义\\\midrule\endhead\bottomrule\endfoot
基线与递减 & 初次响应$\geq5$；$E>0$时计算D；递减判据为$D\geq0.30$。\\
平台 & 末两组三次响应均值差$\leq0.10E$；24次训练时比较19--21与22--24次。\\
首轮恢复 & 已达到递减判据，10s探测$-L\geq0.20E$，且探测$\geq0.80I$。\\
状态干预 & 资源恢复后的响应$\geq0.80I$；重置快状态与sham之差$\leq1$；同时恢复两类状态及旁路比较均要求全脑事件一致。\\
代理预测 & 每条件MAE$\leq1.5$，另比较两个判别条件的MAE与RMSE。\\
H3，10s & 后续初次$\geq0.80I_1$；第三轮D较首轮增加$\geq0.10$，或按本轮初次归一化的连续三次半响应起点提前至少1次，中间轮不反向。\\
H3，2s & 后续初次高于前轮末期；第三轮前3次均值较首轮下降$\geq1$，或按共同首轮基准计算的半响应起点提前至少1次，中间轮不反向。\\
H4(b)主F判据 & 两组均满足基线、递减和平台条件；高频F50区间明确更早，且0.5/1/2s的F优势中位数$\geq0.10$。另报告初次80\%越阈区间。\\
H5 & 两组基线可用，弱输入D较强输入高$\geq0.10$。\\
H6 & 两种训练长度均满足比较条件；24次的半恢复区间更晚，或至少两个探测点比12次少$\geq1$脉冲。\\
H7 & B初次$\geq5$且为A的0.5--2倍；训练后A减弱而B保留至少80\%响应。\\
H8 & B产生至少100个非输入脉冲，B后空白aBN1为0；A相对等时休息的增益$\geq\max(2,0.20I)$，并检查未训练背景增益。\\
一致性 & 补测至少2/3实现符合相应判据，作为描述性一致性标准。\\
\end{longtable}\endgroup

\section{实现检查与运行记录}
各阶段分析程序曾独立回读原始事件，核对源文件和数据哈希、条件覆盖、整数ID、时间网格、响应计数、资源范围及实际感觉输入。$U=0$扩展与M0训练事件一致；训练末态重放、两类状态同时恢复及旁路比较也经过全脑事件核验。H3连续2s恢复的首次探测与独立2s分支一致，H4(b)补测的初始事件与原主程序一致。记录为静默的区间未检出脉冲，期间神经状态仍按模型连续演化。

首轮运行因1800s工具超时中断，随后按原参数分批完成，未完成片段排除。时间网格校验改用严格绝对容差，正式刺激参数未改变。阶段二旁路发生器在恢复空快照时出错，显式清空未使用事件后完成检查及正式运行。补测分析器曾缺少时区导入，修复后重新读取原数据，未重跑神经仿真。不完整的条件组合由分析器拒绝，未生成部分汇总。

首轮有效检查点和批次记录约51.1分钟，阶段二成功批次24.65分钟，补测56.58分钟。这些时间包含数据加载、快照、计算和写盘；失败启动、部分中断开销及排版时间未完整计入，不作为精确总成本或纯积分性能指标。

仿真使用独立环境，主要版本为Python3.10.21、Brian2 2.5.1、NumPy1.22.3、Pandas1.4.3、Cython0.29.33和PyArrow17.0.0，采用单进程执行。PyArrow11退出故障及版本更换记录保留在原工程文档中。实验材料首次公开前，共10+6+16项、合计32项相关单元测试通过。本次文字修订未更改仿真程序、分析算法或依赖。

\section{数据文件与复现说明}
\begin{table}[!htbp]\centering\small
\caption{主要数据与程序路径，相对于\texttt{habituation/}。原始事件目录未随公开仓库提供。}
\begin{tabularx}{\linewidth}{L{3.0cm}Y}\toprule
用途 & 文件或目录\\\midrule
首轮数据与判据 & \path{analysis/metrics_by_seed.csv}、\path{analysis/report.json}、\path{protocols/main_v1.json}\\
状态干预 & \path{stage2/analysis/cycle1/}\\
冻结模型与预测 & \path{stage2/analysis/cycle2/frozen_models.json}、\path{stage2/analysis/cycle2/prospective_predictions.csv}\\
新输出与误差 & \path{stage2/analysis/cycle3/forecast_outcomes.csv}、\path{scores_by_condition.csv}（同目录）\\
特征判定 & \path{hallmarks/analysis/final/report.json}、\path{screens.csv}（同目录）\\
H4(b)补测 & \path{hallmarks/FREQUENCY_SUPPLEMENT.md}、\path{hallmarks/analysis/full/}中的主分析\\
报告构建与检查 & \path{synthesis/prepare_data.py}、\path{synthesis/build.sh}、\path{synthesis/validate.py}\\\bottomrule
\end{tabularx}\end{table}

公开汇总可用于重建本文的表格与PDF，无需启动全脑仿真。若需独立核验原始脉冲，应取得本地事件文件，或按照各阶段README准备上游和环境，将新实验输出到独立目录。新运行产生的数据应与原冻结预测分别记录。上游模型固定提交为\path{91bdd1e7dcf193f3e7ca5a8933497fcef63b7960}。
\FloatBarrier
